\lesson{8}{Mar 04 2022 Fri (22:30:44)}{Remainder to evaluate polynomial}{Unit 4}

\begin{theorem}[Division Algorithm]
  \label{thrm:division_algorithm}

  Let $P(x)$ be any polynomial and $d(x)$ any non constant polynomial or degree
  less than or equal to $P(x)$.

  Then, there exist unique polynomials $Q(x)$ and $R(x)$ such that:

  \begin{equation}
    \label{eqt:division_algorithm}
    \begin{split}
      P(x) = d(x)Q(x) + R(x) \\
    \end{split}
  .\end{equation}

  where either $R(x) = 0$ or the degree of $R(x)$ is less than the degree of
  $d(x)$.
\end{theorem}

\begin{note}
  \label{nte:remainder_theorem}

  Note that we can divide both sides of the equation above by $d(x) \ne 0$ to
  obtain:

  \begin{equation}
    \label{eqt:division_algorithm_remainder}
    \begin{split}
      \frac{P(x)}{d(x)} = Q(x) + \frac{R(x)}{d(x)}
    \end{split}
  .\end{equation}

  This shows the connection with division: when $P(x)$ is divided by $d(x)$, the
  result is $Q(x)$ and the remainder is $R(x)$.
\end{note}

\begin{remark}
  \label{rmk:remainder_theorem}

  \begin{itemize}
    \label{item:remainder_theorem}
  
    \item In the statement of the division algorithm, the polynomial $d(x)$ must
      be \textbf{"Non constant"}

      A constant polynomial $d(x)$ is one of the form $d(x) = c$ for some value
      $c$. Thus, a non constant polynomial is one that is not of this form, that
      is, one whose degree is at least $1$.

    \item A \textbf{Similar Division Algorithm} exists for integers.
    \item If $R(x) = 0$, then  $d(x)$ divided $P(x)$ or divided evenly into
      $P(x)$, or is a factor of $P(x)$
  \end{itemize}
\end{remark}

\begin{theorem}[Remainder Theorem]
  \label{thrm:remainder_theorem}

  If a number $c$ is substituted for $x$ in the polynomial $P(x)$, then the
  result, $P(c)$ is the remainder obtained when $P(x)$ is divided by $x - c$.
\end{theorem}

\begin{proof}
  \label{prf:remainder_theorem}

  By the division algorithm (\ref{thrm:division_algorithm}), when we divided
  $P(x)$ by $x - c$, we get a quotient $Q(x)$ and a remainder $R(x)$. We can
  write this as follows:

  \begin{equation}
    \label{eqt:prf_remainder_theorem_1}
    \begin{split}
      P(x) = (x - c)Q(x) + R(x) \\
    \end{split}
  .\end{equation}

  The degree of $R(x)$ must be less than the degree of $x - c$, so $R(x)$ must
  be constant. This means that $R(x) = r$ for some number $r$.

  \begin{equation}
    \label{eqt:prf_remainder_theorem_2}
    \begin{split}
      P(x) = (x - c)Q(x) + r \\
    \end{split}
  .\end{equation}

  We see that $P(c) = r$:

  \begin{equation}
    \label{eqt:prf_remainder_theorem_3}
    \begin{split}
      P(c)  &= (c - c)Q(c) + r \\
            &= 0 \times Q(c) + r \\
            &= r
    \end{split}
  .\end{equation}

  Thus, $P(c)$ is equal to the remainder when $P(x)$ is divided by $x - c$.
\end{proof}

\subsubsection{Synthetic Division}
\label{sub_sub_sec:synthetic_division}

I'll illustrate these steps by using them to divide
$P(x) = 4x^{2} - x^{4} - 7 - 2x$ by $x + 1$

\begin{enumerate}
  \label{enum:steps_for_synthetic_division}

  \item Arrange the coefficients of $P(x)$ in order of descending powers of $x$.
    Write $0$ as the coefficient for each missing power of $x$.

    \begin{equation}
      \label{eqt:steps_for_synthetic_division}
      \begin{split}
        p(x) = -x^{4} + 0x^{3} - 4x^{2} - 3x - 7 \\
      \end{split}
    .\end{equation}

  \item Insert $r$ (which is any number) to the left of the division symbol.

    Then, bring the leading coefficient of the dividend down.

    \polyhornerscheme[x=-1,stage=2]{-4x^4+0x^3-4x^2-3x-7}

  \item Generate the second and third rows of numbers as follows.

    \begin{enumerate}
      \label{enum:second_and_third_rows}
    
      \item After bringing down the first coefficient of the dividend, multiply
        it by $r$ and place the resulting product under the second coefficient.

        \polyhornerscheme[x=-1,stage=3]{-4x^4+0x^3-4x^2-3x-7}

      \item Add the new number to the second coefficient, and place the sum
        below.

        \polyhornerscheme[x=-1,stage=4]{-4x^4+0x^3-4x^2-3x-7}

      \item Repeat until you've completed all of the coefficients. Here's how it
        should look like after you've completed all steps properly:

        \polyhornerscheme[x=-1]{-4x^4+0x^3-4x^2-3x-7}

    \end{enumerate}

    \item The final entry in the last row is the remainder $R$, and the other
      entries in the row are the coefficients (in descending order) of a
      polynomial whose degree is $1$ less than that of the dividend; This
      polynomial is the quotient $Q(x)$ 

      \begin{equation}
        \label{eqt:final_results}
        \begin{split}
          Q(x)  &= -x^{3} + x^{2} + 3x - 5 \\
          R     &= -2
        \end{split}
      .\end{equation}
\end{enumerate}

% subsubsection synthetic_division (end)

\begin{example}
  \label{exm:remainder_theorem}

  Use the remainder theorem (\ref{thrm:remainder_theorem}) to find $P(2)$ for
  $P(x) = 2x^{3} - 4x^{2} - 9$.
  
  Specifically, give the quotient and the remainder for the associated division
  and the values of $P(2)$.

  So, if $P(x) = 2x^{3} - 4x^{2} - 9$ is divided by $x - 2$, then the remainder
  is $P(2)$.

  We first do the division, rewriting  $P(x) = 2x^{3} - 4x^{2} - 9$ as
  $P(x) = 2x^{3} - 4x^{2} + 0x - 9$:

  \begin{center}
    \polyhornerscheme[x=-2]{2x^3-4x^2+0x-9}
  \end{center}

  So, here's our final result:

  \begin{equation}
    \label{eqt:final_result_2}
    \begin{split}
      Q(x)  &= 2x^{3} - 8x^{2} + 16x \\
      R     &= -41
    \end{split}
  .\end{equation}
\end{example}

\newpage
